% fig-ch0-wonham-supervisory.tex -- redrawn to the house figure language (v2).
% Reader question: which events can the supervisor disable, and what remains
% of the plant's behaviour?
% Claim: the supervisor S restricts the plant G's uncontrolled language
% L(G) to the safe sublanguage K, so L(S/G) subseteq K subset L(G).
% Controllable events Sigma_c (file writes, process spawns, network
% requests) that would leave K are disabled by S; uncontrollable events
% Sigma_u (model token generation, host signals, process exit) always
% execute -- S has no authority over them.
% Grammar chosen: nested sets (tpl-nested-sets.tex) for L(G) superset K,
% beside a small automaton for the plant, with S drawn as a pd trust
% boundary around the controllable-event transitions only; the one
% uncontrollable transition sits outside that boundary, undisabled.
% Rejected: the four-voting-robots-and-a-striped-barrier illustration this
% fragment used to be, plus a processor-chip "kernel shield" icon and six
% small pictograms for individual event kinds (a branching-diamond for git
% commits, a globe for network egress, a gear for OS signals, a stopwatch,
% a lightning bolt) -- none of those drawings states which events are
% controllable, which figure is the actual claim; a set-containment diagram
% plus a labelled automaton states it directly.
% Facts: whitepaper/chapter-0-prerequisites.tex Sec. "Lifecycle Hooks and
% Supervisory Control" (the Ramadge-Wonham formalization, the Sigma_c/Sigma_u
% partition and the L(S/G) subseteq K subset L(G) containment) and its
% caption for fig:wonham-supervisory.
\pdfigurehue{pdcobalt}%
\begin{tikzpicture}[pd figure,x=1cm,y=1cm]
  % ---- one grid, declared once --------------------------------------------
  \def\R{13.90}      % legend rule width
  % ---- left: the language containment, as nested sets ---------------------
  \draw[pd neutral fill] (0.60,1.00) rectangle (6.00,5.30);
  \node[pd label,anchor=north west] at (0.75,5.15) {$L(G)$};
  \draw[pd focus fill] (1.70,1.60) rectangle (4.90,3.90);
  \node[pd label,anchor=north west] at (1.85,3.75) {$K = L(S/G)$};
  \node[pd label] at (3.30,0.65) {$L(S/G) \subseteq K \subset L(G)$};
  \node[pd focus datum] (AD) at (3.00,2.60) {};
  \node[pd badge] at (3.35,2.90) {1};
  \node[pd breach datum] (DD) at (5.35,2.20) {};
  \node[pd badge] at (5.35,2.65) {2};
  % ---- right: the plant, with S as a boundary around Sigma_c only --------
  \draw[pd trust boundary] (7.70,4.60) rectangle (13.00,6.50);
  \node[pd kind,anchor=north] at (10.35,6.44) {S disables $\Sigma_c$};
  \node[pd automaton state] (X0) at (8.60,3.30) {$x_0$};
  \node[pd automaton state] (X1) at (8.60,5.60) {$x_1$};
  \node[pd automaton state,draw=pderror,dash pattern=on 3.2pt off 2.2pt,line width=1.2pt]
    (X3) at (12.40,5.60) {$x_3$};
  \node[pd automaton state] (X2) at (8.60,1.00) {$x_2$};
  \draw[pd initial] (7.30,2.60) -- (X0);
  \draw[pd rule] (X0) -- (X1);
  \node[pd badge] at (8.10,4.45) {1};
  \node[pd label,anchor=east] at (7.75,4.45) {file write};
  \draw[pd breach arrow] (X1) -- (X3);
  \node[pd badge] at (9.85,4.95) {2};
  \node[pd label,anchor=west,text=pderror] at (10.20,4.95) {process spawn};
  \draw[pd rule] (X0) -- (X2);
  \node[pd label,anchor=east] at (8.15,2.15) {model tokens};
  \node[pd note,anchor=north west,text width=6.3cm] at (7.70,0.60)
    {$x_3$ is reached only inside $S$'s boundary; $x_2$, reached by an
    uncontrollable event, is never a candidate for disablement.};
  % ---- the legend: each event class, badge, what S may do -----------------
  \def\ly{-0.90}
  \draw[pd rule] (0,\ly) -- (\R,\ly);
  \node[pd title,anchor=north west] at (.05,{\ly-.30}) {event};
  \node[pd title,anchor=north west] at (3.55,{\ly-.30}) {class};
  \node[pd title,anchor=north west] at (5.35,{\ly-.30}) {the supervisor};
  \draw[pd hairline] (0,{\ly-.95}) -- (\R,{\ly-.95});
  \foreach \i/\n/\h/\l/\s in {%
    0/1/{file write}/{$\Sigma_c$}/{admitted: the trace stays inside $K$},
    1/2/{process spawn}/{$\Sigma_c$}/{disabled: would leave $K$ (e.g.\ an unleased mutation)},
    2/{--}/{model tokens, host signals,\\process exit}/{$\Sigma_u$}/{always executes; outside $S$'s authority}}{
      \node[pd badge] at (.28,{\ly-1.35-\i*.95}) {\n};
      \node[pd label,anchor=base west,align=left,text width=2.75cm] at (.62,{\ly-1.35-\i*.95}) {\h};
      \node[pd label,anchor=base west] at (3.55,{\ly-1.35-\i*.95}) {\l};
      \node[pd label,anchor=base west,align=left,text width=7.9cm] at (5.35,{\ly-1.35-\i*.95}) {\s};}
  \node[pd note,anchor=north west,text width=13.4cm] at (.05,{\ly-1.35-2*.95-1.05})
    {the dot and diamond on the left are the same two transitions, badged
    1 and 2, read as traces inside or outside $K$.};
\end{tikzpicture}
